\documentclass[a4paper]{article}
\usepackage{ctex}
\usepackage[affil-it]{authblk}
\usepackage[backend=bibtex,style=numeric]{biblatex}
\usepackage{amsmath,amsthm,amssymb,amsfonts}
\usepackage{graphicx}
\usepackage{caption, subcaption}
\usepackage{bm}

\usepackage{geometry}
\geometry{margin=1.5cm, vmargin={0pt,1cm}}
\setlength{\topmargin}{-1cm}
\setlength{\paperheight}{29.7cm}
\setlength{\textheight}{25.3cm}

\addbibresource{citation.bib}

\begin{document}
% =================================================
\title{微分方程数值解 2026 春夏作业六}

\author{曾申昊 3220100701
  \thanks{邮箱: \texttt{3097714673@qq.com}
                            \texttt{或} \texttt{3220100701@zju.edu.cn}}}
\affil{浙江大学}

\date{更新时间: \today}

\maketitle

% ============================================

\section*{Exercise 13.28}

\par 三种方法的 Butcher 表如下所示。

\begin{table}[htbp]
    \centering
    \caption{Butcher tableaux of different methods}
    \begin{subtable}[t]{0.3\linewidth}
        \centering
        \caption{The modified Euler method}
        \begin{tabular}{c|cc}
        0 & 0 & ~ \\
        $\frac{1}{2}$ & $\frac{1}{2}$ & 0 \\
        \hline
        ~ & 0 & 1 \\
    \end{tabular}
    \end{subtable}
    \begin{subtable}[t]{0.3\linewidth}
        \centering
        \caption{The improved Euler method}
        \begin{tabular}{c|cc}
            0 & 0 & ~ \\
            1 & 1 & 0 \\
            \hline
            ~ & $\frac{1}{2}$ & $\frac{1}{2}$ \\
        \end{tabular}
    \end{subtable}
    \begin{subtable}[t]{0.3\linewidth}
        \centering
        \caption{The Heun's third-order method}
        \begin{tabular}{c|ccc}
            0 & 0 & ~ & ~\\
            $\frac{1}{3}$ & $\frac{1}{3}$ & 0 & ~\\
            $\frac{2}{3}$ & 0 & $\frac{2}{3}$ & 0\\
            \hline
            ~ & $\frac{1}{4}$ & 0 & $\frac{3}{4}$\\
        \end{tabular}
    \end{subtable}
\end{table}

\section*{Exercise 13.29}

\par 由 \textbf{Definition 13.26} 可知，
一个 s 步的 Runge-Kutta 方法的通用形式为
\begin{equation}
    \left\{
        \begin{aligned}
            &\mathbf{y}_i = \mathbf{f}\left(
                \mathbf{U}^n + k \sum_{j=1}^{s} a_{i,j} \mathbf{y}_j, t_n + c_i k
            \right)\\
            &\mathbf{U}^{n+1} = \mathbf{U}^n + k \sum_{j=1}^{s} b_j \mathbf{y}_j
        \end{aligned}
    \right.
\end{equation}
\par 我们令
$\bm{\xi}_i=\mathbf{U}^n + k \sum_{j=1}^{s} a_{i,j} \mathbf{y}_j$，
代入第一式，得到
\begin{equation}
    \label{eq:xi to y}
    \mathbf{y}_i = \mathbf{f}(\bm{\xi}_i, t_n + c_i k)
\end{equation}
\par 更换指标 $i$ 为 $j$，两边同乘 $a_{i,j}$ 和 $k$，
并对 $j$ 从 $1$ 到 $s$ 求和，最后加上 $\mathbf{U}^n$ 得到
\begin{equation}
    \mathbf{U}^n + k \sum_{j=1}^{s} a_{i,j} \mathbf{y}_j = \mathbf{U}^n + k \sum_{j=1}^{s} a_{i,j} \mathbf{f}(\bm{\xi}_j, t_n + c_j k)
\end{equation}
\par 也就是
\begin{equation}
    \bm{\xi}_i = \mathbf{U}^n + k \sum_{j=1}^{s} a_{i,j} \mathbf{f}(\bm{\xi}_j, t_n + c_j k)
\end{equation}
\par 将式 (\ref{eq:xi to y}) 代入第二式得到
\begin{equation}
    \mathbf{U}^{n+1} = \mathbf{U}^n + k \sum_{j=1}^{s} b_j \mathbf{f}(\bm{\xi}_j, t_n + c_j k)
\end{equation}

\section*{Exercise 13.36}

\par 对于三步的 ERK 方法的通式为
\begin{equation}
    \left\{
        \begin{aligned}
            &\mathbf{y}_1 = \mathbf{f}(\mathbf{U}^n, t_n)\\
            &\mathbf{y}_2 = \mathbf{f}(\mathbf{U}^n + k a_{2,1} \mathbf{y}_1, t_n + c_2 k)\\
            &\mathbf{y}_3 = \mathbf{f}(\mathbf{U}^n + k a_{3,1} \mathbf{y}_1 + k a_{3,2} \mathbf{y}_2, t_n + c_3 k)\\
            &\mathbf{U}^{n+1} = \mathbf{U}^n + k b_1 \mathbf{y}_1 + k b_2 \mathbf{y}_2 + k b_3 \mathbf{y}_3
        \end{aligned}
    \right.
\end{equation}
\par 一步误差 $\mathcal{L}\mathbf{u}(t_n)$ 为
\begin{equation}
    \begin{aligned}
        \mathcal{L}\mathbf{u}(t_n)
        =&\ \mathbf{u}(t_{n+1})-\mathbf{u}(t_n)
        -k (b_1 \mathbf{y}_1 + b_2 \mathbf{y}_2+b_3 \mathbf{y}_3)\\
        =&\ \left(
            \mathbf{u}+k\mathbf{u}'+\frac{k^2}{2}\mathbf{u}''+\frac{k^3}{6}\mathbf{u}'''
            +O(k^4)
        \right)-\mathbf{u}
        -k b_1\mathbf{u}'\\
        &\ -k b_2\left(
            \mathbf{f}+k a_{2,1}\mathbf{f}_{\mathbf{u}}\mathbf{y}_1+k c_2\mathbf{f}_t
            +\frac{k^2}{2}a_{2,1}^2 \mathbf{f}_{\mathbf{uu}}\mathbf{y}_1^2
            +\frac{k^2}{2} c_2^2 \mathbf{f}_{tt}
            +k^2 a_{2,1} c_2 \mathbf{f}_{\mathbf{u}t}\mathbf{y}_1
            + O(k^3)
        \right)\\
        &\ -k b_3\left(
            \mathbf{f}
            +k \mathbf{f}_{\mathbf{u}}\big(a_{3,1}\mathbf{y}_1
            +a_{3,2}(\mathbf{f}+k a_{2,1}\mathbf{f}_{\mathbf{u}}\mathbf{y}_1+k c_2\mathbf{f}_t)\big)+k c_3\mathbf{f}_t
            +\frac{k^2}{2}\mathbf{f}_{\mathbf{uu}}(a_{3,1}\mathbf{y}_1
            +a_{3,2}\mathbf{f})^2
            \right.\\
            &\qquad\qquad \left.
            +\frac{k^2}{2} c_3^2 \mathbf{f}_{tt}
            +k^2 c_3 \mathbf{f}_{\mathbf{u}t}(a_{3,1}\mathbf{y}_1
            +a_{3,2}\mathbf{f})
            + O(k^3)
        \right)\\
        =&\ \left(
            \mathbf{u}+k\mathbf{u}'+\frac{k^2}{2}\mathbf{u}''+\frac{k^3}{6}\mathbf{u}'''
            +O(k^4)
        \right)-\mathbf{u}
        -k b_1\mathbf{u}'\\
        &\ -k b_2\left(
            \mathbf{u}'+k c_2(\mathbf{f}_{\mathbf{u}}\mathbf{f}+\mathbf{f}_t)
            +\frac{k^2}{2}c_2^2 (\mathbf{f}_{\mathbf{uu}}\mathbf{f}^2
            +\mathbf{f}_{tt}
            +2 \mathbf{f}_{\mathbf{u}t}\mathbf{f})
            + O(k^3)
        \right)\\
        &\ -k b_3\left(
            \mathbf{u}'
            +k c_3(\mathbf{f}_{\mathbf{u}}\mathbf{f}+\mathbf{f}_t)
            +k^2 a_{3,2} c_2 \mathbf{f}_{\mathbf{u}}(
                \mathbf{f}_{\mathbf{u}}\mathbf{f}+\mathbf{f}_t
            )
            +\frac{k^2}{2} c_3^2 (
                \mathbf{f}_{\mathbf{uu}}\mathbf{f}^2
                +\mathbf{f}_{tt}
                +2 \mathbf{f}_{\mathbf{u}t}\mathbf{f}
            )
            + O(k^3)
        \right)\\
        =&\ \left(
            \mathbf{u}+k\mathbf{u}'+\frac{k^2}{2}\mathbf{u}''+\frac{k^3}{6}\mathbf{u}'''
            +O(k^4)
        \right)-\mathbf{u}
        -k b_1\mathbf{u}'\\
        &\ -k b_2\left(
            \mathbf{u}'+k c_2\mathbf{u}''
            +\frac{k^2}{2}c_2^2 (\mathbf{u}'''-\mathbf{f}_{\mathbf{u}}\mathbf{u}'')
            + O(k^3)
        \right)\\
        &\ -k b_3\left(
            \mathbf{u}'
            +k c_3\mathbf{u}''
            +k^2 a_{3,2} c_2 \mathbf{f}_{\mathbf{u}}\mathbf{u}''
            +\frac{k^2}{2} c_3^2 (\mathbf{u}'''-\mathbf{f}_{\mathbf{u}}\mathbf{u}'')
            + O(k^3)
        \right)\\
        =&\ k\left(
            1 - b_1 - b_2 - b_3
        \right)\mathbf{u}'
        +k^2\left(
            \frac{1}{2} - b_2 c_2 - b_3 c_3
        \right)\mathbf{u}''\\
        &\ +k^3\left(
            \frac{1}{6} - \frac{1}{2}b_2c_2^2 - \frac{1}{2}b_3 c_3^2
        \right)\mathbf{u}'''
        + k^3 \left(
            \frac{1}{2}b_2 c_2^2+\frac{1}{2}b_3 c_3^2 - a_{3,2}b_3c_2
        \right)\mathbf{f}_{\mathbf{u}}\mathbf{u}''
        + O(k^4)
    \end{aligned}
\end{equation}
\par 可以看出该方法的阶数为 3 的充要条件为
\begin{equation}
    \left\{
        \begin{aligned}
            &b_1 + b_2 + b_3 = 1\\
            &b_2 c_2 + b_3 c_3 = \frac{1}{2}\\
            &b_2 c_2^2 + b_3 c_3^2 = \frac{1}{3}\\
            &a_{3,2}b_3c_2 = \frac{1}{6}
        \end{aligned}
    \right.
\end{equation}
\par 因此可以验证下面的 Butcher 表是满足该方程的，
但是 Heun's third-order formula 并不属于该方法类。
\begin{table}[htbp]
    \centering
    \caption{Butcher tableau of the three-step ERK method}
    \begin{tabular}{c|ccc}
        0 & 0 & ~ & ~\\
        $\frac{2}{3}$ & $\frac{2}{3}$ & 0 & ~\\
        $\frac{2}{3}$ & $\frac{2}{3}-\frac{1}{4\alpha}$ & $\frac{1}{4\alpha}$ & 0\\
        \hline
        ~ & $\frac{1}{4}$ & $\frac{3}{4}-\alpha$ & $\alpha$\\
    \end{tabular}
\end{table}

\section*{Exercise 13.42}

\par “$\Rightarrow$” 的证明：
RK 方法对于 $\mathbb{P}_{r-1}$ 中的任意多项式的数值积分都是精确的，
不妨取 $f=(x - t_n)^{l-1}$，其中 $l = 1,2,\ldots,r$。由数值积分的定义可知
\begin{equation}
    \begin{aligned}
        I_s((x-t_n)^{l-1})
        &=k\sum_{j=1}^{s} b_j \left(c_j k \right)^{l-1}\\
        &=k^l \sum_{j=1}^{s} b_j c_j^{l-1}
    \end{aligned}
\end{equation}
\par 另一方面，$(x-t_n)^{l-1}$ 的积分为
\begin{equation}
    \int_{t_n}^{t_n+k} (x-t_n)^{l-1} \text{d}x = \frac{1}{l} k^l
\end{equation}
\par 对比可得 $\displaystyle\sum_{j=1}^{s} b_j c_j^{l-1} = \dfrac{1}{l}$，
也就是说 RK 方法是 $B(r)$ 的。

\par “$\Leftarrow$” 的证明：RK 方法是 $B(r)$ 的，也就是对于
$\forall l = 1,2,\ldots,r$ 有 
$\displaystyle\sum_{j=1}^{s} b_j c_j^{l-1} = \dfrac{1}{l}$。
下面计算 $x^{l-1}$ 的数值积分
\begin{equation}
    \begin{aligned}
        I_s(x^{l-1})&=
        k\sum_{j=1}^{s} b_j \left( t_n + c_j k \right)^{l-1}\\
        &=
        k\sum_{j=1}^{s} b_j \sum_{m=0}^{l-1} \binom{l-1}{m} t_n^{l-1-m} c_j^m k^m\\
        &=
        k\sum_{m=0}^{l-1} \binom{l-1}{m} t_n^{l-1-m} k^m \sum_{j=1}^{s} b_j c_j^m\\
        &=
        k\sum_{m=0}^{l-1} \binom{l-1}{m} t_n^{l-1-m} k^m \frac{1}{m+1}\\
        &=\frac{1}{l}\sum_{m=0}^{l-1} \binom{l}{m+1} t_n^{l-1-m} k^{m+1}\\
        &=\frac{1}{l}\left[
            (t_n+k)^l - t_n^l
        \right]
        =\int_{t_n}^{t_{n}+k} x^{l-1} \text{d}x
    \end{aligned}
\end{equation}
\par 也就是说 RK 方法对于 $x^{l-1}$ 的数值积分是精确的。由
线性组合可得 RK 方法对于 $\mathbb{P}_{r-1}$ 中的任意多项式的数值积分都是精确的。

\section*{Exercise 13.59}

\par 假设精确解为 $\mathbf{u}(t)$，在区间
$[t_n, t_{n+1}]$ 构造多项式 $\mathbf{q}(t)$ 满足
\begin{equation}
    \left\{
        \begin{aligned}
            &\mathbf{q}(t_n) = \mathbf{u}(t_n)\\
            &\mathbf{q}'(t_n^{(i)}) = \mathbf{u}'(t_n^{(i)})
            =\mathbf{f}\left(
                \mathbf{u}(t_n^{(i)}),t_n^{(i)}
            \right)
            \qquad i = 1,2,\ldots,s
        \end{aligned}
    \right.
\end{equation}
\par 为了证明该方法的阶数至少为 $s$，
我们需要证明
\begin{equation}
    \mathcal{L}\mathbf{u}(t_n) 
    = \mathbf{u}(t_{n+1}) - \mathbf{p}(t_{n+1}) = \left(\mathbf{u}(t_{n+1}) - \mathbf{q}(t_{n+1})\right)
    +\left(\mathbf{q}(t_{n+1}) - \mathbf{p}(t_{n+1})\right)
    = O(k^{s+1})
\end{equation}
\par 接下来我们分别证明 $\left\|\mathbf{u}(t) - \mathbf{q}(t)\right\| = O(k^{s+1})$ 和
$\left\|\mathbf{q}(t) - \mathbf{p}(t)\right\| = O(k^{s+1})$。
\begin{itemize}
    \item $\left\|\mathbf{u}(t) - \mathbf{q}(t)\right\| = O(k^{s+1})$\\
    $\mathbf{p}(t)$ 是 $\mathbf{u}(t)$ 的插值函数，满足1个点的函数值和
    $s$ 个点的导数值相等。这一部分由插值误差的估计可得。
    \item $\left\|\mathbf{q}(t) - \mathbf{p}(t)\right\| = O(k^{s+1})$\\
    由 \textbf{Definition 13.56} 可知
    \begin{equation}
        \left\{
            \begin{aligned}
                &\mathbf{p}(t_n) = \mathbf{u}(t_n)\\
                &\mathbf{p}'(t_n^{(i)}) = \mathbf{f}\left(
                    \mathbf{p}(t_n^{(i)}),t_n^{(i)}
                \right)
                \qquad i = 1,2,\ldots,s
            \end{aligned}
        \right.
    \end{equation}
    于是我们有
    \begin{equation}
        \begin{gathered}
            \left\|\mathbf{q}'(t_n^{(i)}) - \mathbf{p}'(t_n^{(i)})\right\|
            = \left\|\mathbf{f}\left(
                \mathbf{u}(t_n^{(i)}),t_n^{(i)}
            \right) - \mathbf{f}\left(
                \mathbf{p}(t_n^{(i)}),t_n^{(i)}
            \right)\right\|\\
            \left\|\mathbf{q}'(t_n^{(i)}) - \mathbf{p}'(t_n^{(i)})\right\|
            \leq L\left\|
                \mathbf{u}(t_n^{(i)}) - \mathbf{p}(t_n^{(i)})
            \right\|\\
            \left\|\mathbf{q}'(t_n^{(i)}) - \mathbf{p}'(t_n^{(i)})\right\|
            \leq L\left\|
                \mathbf{q}(t_n^{(i)}) - \mathbf{p}(t_n^{(i)})
            \right\| + O(k^{s+1})\\
        \end{gathered}
    \end{equation}
    由于 $\mathbf{q}(t)$ 和 $\mathbf{p}(t)$ 都是多项式，
    所以可以表示为
    \begin{equation}
        \mathbf{q}(t) - \mathbf{p}(t) 
        = \sum_{j=1}^{s} \left(
            \mathbf{q}'(t_n^{(i)}) - \mathbf{p}'(t_n^{(i)})
        \right)\int_{t_n}^{t} \ell_j(\tau) \text{d}\tau
    \end{equation}
    其中 $\ell_j(\tau)$ 是 Lagrange 基函数。于是我们有
    \begin{equation}
        \begin{aligned}
            \left\|\mathbf{q}(t) - \mathbf{p}(t)\right\|
            &\leq \sum_{j=1}^{s} \left\|
                \mathbf{q}'(t_n^{(i)}) - \mathbf{p}'(t_n^{(i)})
            \right\| \int_{t_n}^{t} |\ell_j(\tau)| \text{d}\tau\\
            &\leq k\sum_{j=1}^{s} \left\|
                \mathbf{q}'(t_n^{(i)}) - \mathbf{p}'(t_n^{(i)})
            \right\| \int_{0}^{c} |\ell_j(\xi)| \text{d}\xi\\
            &\leq Mk\sum_{j=1}^{s} \left\|
                \mathbf{q}'(t_n^{(i)}) - \mathbf{p}'(t_n^{(i)})
            \right\|\\
        \end{aligned}
    \end{equation}
    第二行做了变量代换 $\xi = \frac{\tau - t_n}{k}$，使得积分与 $k$ 无关；
    第三行中 $M = \displaystyle\max_{j=1,2,\ldots,s} \int_{t_n}^{t} |\ell_j(\tau)| \text{d}\tau$。
    做记号 $E = \displaystyle\max_{j=1,2,\ldots,s} \left\|
        \mathbf{q}(t_n^{(i)}) - \mathbf{p}(t_n^{(i)})
    \right\|$ 和 $E' = \displaystyle\max_{j=1,2,\ldots,s} \left\|
        \mathbf{q}'(t_n^{(i)}) - \mathbf{p}'(t_n^{(i)})
    \right\|$，则上两式子可以写为
    \begin{equation}
        \begin{cases}
            E' \leq L E + O(k^{s+1})\\
            E \leq Mks E'
        \end{cases}
    \end{equation}
    不难得到 $E' = O(k^{s+1})$ 和 $E = O(k^{s+2})$，
    最后易证 $\left\|\mathbf{q}(t) - \mathbf{p}(t)\right\| 
    \leq Mk\displaystyle\sum_{j=1}^{s} \left\|
        \mathbf{q}'(t_n^{(i)}) - \mathbf{p}'(t_n^{(i)})
    \right\| = O(k^{s+1})$。
\end{itemize}

\section*{Exercise 13.60}

\par 由 \textbf{Theorem 13.58} 可知，
$a_{i,j} = \displaystyle\int_{0}^{c_i}\ell_j(\tau) \text{d}\tau$
和 $b_j = \displaystyle\int_{0}^{1}\ell_j(\tau) \text{d}\tau$。
\par 下面计算 $a_{i,j}$ 的求和
\begin{equation}
    \begin{aligned}
        \sum_{j=1}^{s} a_{i,j} &= \sum_{j=1}^{s} \int_{0}^{c_i}\ell_j(\tau) \text{d}\tau\\
        &= \int_{0}^{c_i} \sum_{j=1}^{s} \ell_j(\tau) \text{d}\tau\\
        &= \int_{0}^{c_i} 1 \text{d}\tau = c_i
    \end{aligned}
\end{equation}
\par 同理可得 $\displaystyle\sum_{j=1}^{s} b_j = 1$。

\section*{Exercise 13.62}

\par 取 $c_1=\frac{1}{4}$，$c_2=\frac{1}{2}$，$c_3=\frac{3}{4}$，对应的
Lagrange 基函数为
\begin{equation}
    \begin{aligned}
        \ell_1(\tau) &= \frac{(\tau-\frac{1}{2})(\tau-\frac{3}{4})}{(\frac{1}{4}-\frac{1}{2})(\frac{1}{4}-\frac{3}{4})} = 8\tau^2 - 10\tau + 3\\
        \ell_2(\tau) &= \frac{(\tau-\frac{1}{4})(\tau-\frac{3}{4})}{(\frac{1}{2}-\frac{1}{4})(\frac{1}{2}-\frac{3}{4})} = -16\tau^2 + 16\tau - 3\\
        \ell_3(\tau) &= \frac{(\tau-\frac{1}{4})(\tau-\frac{1}{2})}{(\frac{3}{4}-\frac{1}{4})(\frac{3}{4}-\frac{1}{2})} = 8\tau^2 - 6\tau + 1
    \end{aligned}
\end{equation}
\par 对应的积分为
\begin{equation}
    \begin{aligned}
        \int_0^t \ell_1(\tau) \text{d}\tau &= \frac{8}{3}t^3 - 5t^2 + 3t\\
        \int_0^t \ell_2(\tau) \text{d}\tau &= -\frac{16}{3}t^3 + 8t^2 - 3t\\
        \int_0^t \ell_3(\tau) \text{d}\tau &= \frac{8}{3}t^3 - 3t^2 + t
    \end{aligned}
\end{equation}
\par 然后可以计算三步 IRK 方法的 Butcher 表
\begin{table}[htbp]
    \centering
    \caption{Butcher tableau of the three-step IRK method}
    \begin{tabular}{c|ccc}
        $\frac{1}{4}$ & $\frac{23}{48}$ & $-\frac{1}{3}$ & $\frac{5}{48}$\\
        $\frac{1}{2}$ & $\frac{7}{12}$ & $-\frac{1}{6}$ & $\frac{1}{12}$\\
        $\frac{3}{4}$ & $\frac{9}{16}$ & 0 & $\frac{3}{16}$\\
        \hline
        ~ & $\frac{2}{3}$ & $-\frac{1}{3}$ & $\frac{2}{3}$\\
    \end{tabular}
\end{table}

\section*{Exercise 13.65}

\par $B(s+r)$ 的条件为对 $\forall l = 1,2,\ldots,s+r$ 有 
\begin{equation}
    \sum_{j=1}^{s} b_j c_j^{l-1} = \frac{1}{l}
\end{equation}
\par $C(s)$ 的条件为对 $\forall i = 1,2,\ldots,s$ 和 $\forall m = 1,2,\ldots,s$ 有
\begin{equation}
    \sum_{j=1}^{s} a_{i,j} c_j^{m-1} = \frac{c_i^m}{m}
\end{equation}
\par 定义向量 $\mathbf{u}$ 和 $\mathbf{v}$ 为
\begin{gather}
    u_j = \sum_{i=1}^{s} b_ic_i^{m-1}a_{i,j}\\
    v_j = \frac{1}{m}b_j(1 - c_j^m)
\end{gather}
\par 分别计算 $V\mathbf{u}$ 和 $V\mathbf{v}$，
其中 $V_{i,j} = c_j^{i-1}$。
\begin{equation}
    \begin{aligned}
        \left(V\mathbf{u}\right)_i
        =&\ \sum_{j=1}^{s} V_{i,j} u_j\\
        =&\ \sum_{j=1}^{s} c_j^{i-1} \sum_{k=1}^{s} b_k c_k^{m-1} a_{k,j}\\
        =&\ \sum_{k=1}^{s} b_k c_k^{m-1} \sum_{j=1}^{s} a_{k,j} c_j^{i-1}\\
        =&\ \frac{1}{i}\sum_{k=1}^{s} b_k c_k^{i+m-1}\\
        =&\ \frac{1}{i(i+m)}
    \end{aligned}
\end{equation}
\par 另一边
\begin{equation}
    \begin{aligned}
        \left(V\mathbf{v}\right)_i
        =&\ \sum_{j=1}^{s} V_{i,j} v_j\\
        =&\ \sum_{j=1}^{s} c_j^{i-1} \frac{1}{m} b_j (1 - c_j^m)\\
        =&\ \frac{1}{m}\sum_{j=1}^{s} b_j c_j^{i-1} - \frac{1}{m}\sum_{j=1}^{s} b_j c_j^{i+m-1}\\
        =&\ \frac{1}{m}\frac{1}{i}-\frac{1}{m}\frac{1}{i+m}\\
        =&\ \frac{1}{i(i+m)}
    \end{aligned}
\end{equation}
\par 也就是说 $V\mathbf{u} = V\mathbf{v}$，由于 $V$ 是可逆矩阵，
所以 $\mathbf{u} = \mathbf{v}$，也就是 $D(r)$ 成立。

\section*{Exercise 13.69}

\begin{equation}
    \mathbf{b} = \left[
        \frac{2}{3}, -\frac{1}{3}, \frac{2}{3}
    \right]^T, \qquad
    \mathbf{c} = \left[
        \frac{1}{4}, \frac{1}{2}, \frac{3}{4}
    \right]^T
\end{equation}
\begin{equation}
    A = \begin{bmatrix}
        \frac{23}{48} & -\frac{1}{3} & \frac{5}{48}\\
        \frac{7}{12} & -\frac{1}{6} & \frac{1}{12}\\
        \frac{9}{16} & 0 & \frac{3}{16}
    \end{bmatrix}, \qquad
    C = \begin{bmatrix}
        \frac{1}{4} & 0 & 0\\
        0 & \frac{1}{2} & 0\\
        0 & 0 & \frac{3}{4}
    \end{bmatrix}
\end{equation}
\par 由 \textbf{Exercise 13.59} 和 \textbf{Corollary 13.70} 
可知，该方法的阶数在 3 到 6 之间。下面先计算 $B(r)$
\begin{equation}
    \begin{aligned}
        \sum_{j=1}^{3} b_j c_j^{1-1}
        =&\ \frac{2}{3}-\frac{1}{3}+\frac{2}{3} = 1\\
        \sum_{j=1}^{3} b_j c_j^{2-1}
        =&\ \frac{2}{3}\cdot\frac{1}{4}-\frac{1}{3}\cdot\frac{1}{2}+\frac{2}{3}\cdot\frac{3}{4} = \frac{1}{2}\\
        \sum_{j=1}^{3} b_j c_j^{3-1}
        =&\ \frac{2}{3}\cdot\left(\frac{1}{4}\right)^2-\frac{1}{3}\cdot\left(\frac{1}{2}\right)^2+\frac{2}{3}\cdot\left(\frac{3}{4}\right)^2 = \frac{1}{3}\\
        \sum_{j=1}^{3} b_j c_j^{4-1}
        =&\ \frac{2}{3}\cdot\left(\frac{1}{4}\right)^3-\frac{1}{3}\cdot\left(\frac{1}{2}\right)^3+\frac{2}{3}\cdot\left(\frac{3}{4}\right)^3 = \frac{1}{4}\\
        \sum_{j=1}^{3} b_j c_j^{5-1}
        =&\ \frac{2}{3}\cdot\left(\frac{1}{4}\right)^4-\frac{1}{3}\cdot\left(\frac{1}{2}\right)^4+\frac{2}{3}\cdot\left(\frac{3}{4}\right)^4 =\frac{74}{384}\neq \frac{1}{5}\\
    \end{aligned}
\end{equation}
\par $B(5)$ 不成立，所以接下来只需验证该方法是否有 4 阶。
相应的要 $(l, m)$ 分别为 $(1,3)$、$(2,2)$、$(3,1)$ 时有
\begin{equation}
    \mathbf{b}^TA^mC^{l-1}\mathbf{1} = \frac{(l-1)!}{(l+m)!}
\end{equation}
\par 具体代入计算得到
\begin{equation}
    \begin{aligned}
        \mathbf{b}^TA^3\mathbf{1} =&\ \frac{1}{24} = \frac{0!}{4!}\\
        \mathbf{b}^TA^2C\mathbf{1} =&\ \frac{1}{24} = \frac{1!}{4!}\\
        \mathbf{b}^TAC^2\mathbf{1} =&\ \frac{1}{12} = \frac{2!}{4!}
    \end{aligned}
\end{equation}
\par 综上所述，该方法的阶数为 4。

% ===============================================

\printbibliography

\end{document}